% ----- Consignes exo 18 ----- %
\begin{td-exo}[\textsc{Vertex Cover} dans les planaires utilisant une 4-coloration]\, % 18 
	% insert algorithm here
	\begin{enumerate}
		\item Rappeler le résultat de complexité pour le problème de l'étape 4-coloration.

		\item Expliquer le fonctionnement de l'algorithme 6. % TODO: add ref to algorithm 6

		\item Prouver la correction de l'algorithme 6. % TODO: add ref to algorithm 6

		\item Montrer que le problème admet un ratio de \(\frac{3}{2}\).
	\end{enumerate}
\end{td-exo}

% ----- Solutions exo 18 ----- %
\iftoggle{showsolutions}{ 
	\begin{td-sol}[]\ % 18
		\begin{enumerate}
			\item Le problème de la 4-coloration dans un graphe planaire est polynomial.

			\item x

			\item La couleur retenue est au moins de taille \(\frac{w}{4}\) où \(w\) est
			\begin{equation*}
				w = \sum_{v \in V} w(v)
			\end{equation*}

			On conserve les sommets de toutes les autres couleurs, ils couvrent donc les sommets de la couleur retenue. 
			Comme ces sommets forment un stable, on couvre bien toutes les arêtes du graphe.

			\item On a
			\begin{equation*}
				OPT_{LP} = \sum_{\substack{v\in V\\x^\ast_v = 1}} x_v c(v) + \sum_{\substack{v\in V\\x^\ast_v = \frac{1}{2}}} x_v c(v) 
			\end{equation*}
			Il en suit
			\begin{equation*}
				\begin{aligned}
					S 
					&= OPT_{LP} + \frac{1}{2} \mathsf{Cost}(\ol k) - \frac{1}{2} \mathsf{Cost}(k)\\
					&\leq OPT_{LP} + \frac{w}{4} \\
					&\leq OPT_{LP} + \frac{1}{2} OPT_{LP} \\
					&\leq \frac{3}{2} OPT_{LP} \\
					&\leq \frac{3}{2} OPT_{ILP} \\
					& \implies \rho \leq \frac{3}{2}
				\end{aligned}
			\end{equation*}
		\end{enumerate}
	\end{td-sol}
}{}
